\documentclass{article}

% --- Control global line spacing for body text ---
\usepackage{setspace}
\setstretch{0.9}     % Multiplier for line spacing; 1 = normal, <1 = tighter lines

% --- Section and subsection heading spacing ---
\usepackage{titlesec}
\titlespacing*{\section}{0pt}{0.1ex}{0pt}
\titlespacing*{\subsection}{0pt}{0.05ex}{0pt}

% --- Section formatting with integrated horizontal line using titlerule ---
\titleformat{\section}
{\normalfont\Large\bfseries}  % format
{}                             % label
{0em}                          % sep between label and title
{}                             % before-code
[\vspace{0.05cm}\titlerule\vspace{0.05cm}]  % after-code: adds line after section title

% --- Make subsection titles larger ---
\titleformat{\subsection}{\normalfont\large\bfseries}{}{0em}{}

% --- Custom command for slightly larger item text ---
\newcommand{\itemsize}{\fontsize{11.5pt}{13.5pt}\selectfont}

% --- Bulleted/numbered list spacing configuration ---
\usepackage{enumitem}
\setlist[itemize]{
	leftmargin=15pt,   % Indentation from the left margin for bullets
	itemsep=2pt,       % Vertical space between list items
	topsep=0pt,        % Space above and below the entire list
	parsep=0pt,        % Space between paragraphs inside an item
	partopsep=0pt      % Extra space added when environment starts a new paragraph
}
\setlist[enumerate]{
	leftmargin=15pt,
	itemsep=2pt,
	topsep=0pt,
	parsep=0pt,
	partopsep=0pt
}

% --- Table column formatting ---
\usepackage{array}

% --- Font configuration ---
\usepackage[quiet]{fontspec}
\setmainfont[
	Ligatures={TeX, Common},             % Enable TeX and common ligatures
	Numbers=Lining,                      % Use lining style numbers
	BoldFont={Linux Libertine O Bold},   % Specify bold font variant
	BoldFeatures={Weight=3}              % Increase bold weight for emphasis
]{Linux Libertine O}

% --- Fix math font error by using unicode-math and Libertinus Math ---
\usepackage{unicode-math}
\setmathfont{Libertinus Math}  % Scalable math font compatible with Linux Libertine

% --- Icon and hyperlink support ---
\usepackage{fontawesome5}              % FontAwesome icons for symbols
\usepackage{hyperref}
\hypersetup{
	colorlinks=true,               % Enable colored links
	linkcolor=black,
	filecolor=black,
	urlcolor=black,
}

% --- Page margin and layout settings ---
\usepackage[a4paper, margin=0.5cm]{geometry} % Narrow margins for dense text

% --- Multi-line comment environment (used for site-only metadata) ---
% Provides \begin{comment}...\end{comment} blocks that are invisible
% in the rendered PDF. The portfolio site parses these for things like
% paragraph project descriptions that would otherwise crowd the resume.
\usepackage{verbatim}

% --- Remove paragraph indent and extra paragraph spacing ---
\setlength{\parindent}{0pt}  % No paragraph indentation
\setlength{\parskip}{0pt}    % No extra space between paragraphs

\begin{document}
\fontsize{11pt}{13pt}\selectfont  % Font size 11pt, baseline skip 13pt (line height) to avoid font errors

% --- Safe & Compact Title Section ---
% -----------------------------------
% HEADER (Compact & Readable)
% -----------------------------------
\begin{center}
	{\Huge \textbf{Rishabh Sahu}} \\ \vspace{3pt}

	% multiset<int>towerTops;Address: Slightly smaller than body text to keep it on one line
	{\fontsize{12pt}{12.5pt}\selectfont \faIcon{map-marker-alt}\enspace \textbf{Address:} Lig 139 Harshwardhan Nagar, Bhopal, Madhya Pradesh, India} \\ \vspace{3pt}

	% Contact Info: 10.5pt Font
	{\fontsize{12pt}{12.5pt}\selectfont
		\faIcon{envelope}\enspace \textbf{Email:} \href{mailto:rishabhsahu617@outlook.com}{rishabhsahu617@outlook.com}
		\quad \textbullet \quad
		\faIcon{mobile-alt}\enspace \textbf{Phone:} (+91) 7898994078
	} \\ \vspace{3pt}

	% Links: 10.5pt Font
	{\fontsize{12pt}{12.5pt}\selectfont
		\faIcon{linkedin}\enspace \textbf{LinkedIn:} \href{https://www.linkedin.com/in/rishabh-sahu325}{rishabh-sahu325}
		\quad \textbullet \quad
		\faIcon{github}\enspace \textbf{GitHub:} \href{https://github.com/RishabhSahuIIIT}{RishabhSahuIIIT}
		\quad \textbullet \quad
		\faIcon{gitlab}\enspace \textbf{GitLab:} \href{https://gitlab.com/rishabhsahu325}{rishabhsahu325}
		\quad \textbullet \quad
		\faIcon{globe}\enspace \textbf{Portfolio:} \href{https://rishabhsahuiiit.github.io/}{rishabhsahuiiit.github.io}
	}
\end{center}
% -----------------------------------
\section*{Education}
\begin{flushleft}
	\noindent\large\textbf{International Institute of Information Technology Hyderabad} \hfill July 2024 -- Present \\
	MTech. Computer Science \hfill CGPA: 7.73/10\\
	\vspace{0.1cm}
	\noindent\large\textbf{The LNM Institute of Information Technology Jaipur} \hfill July 2019 -- June 2023 \\
	BTech. Computer Science \hfill CGPA: 7.43/10\\
	\vspace{0.1cm}
	\noindent\large\textbf{G.V.N. The Global School Bhopal} \hfill April 2017 -- March 2019 \\ High School \hfill 87.2\%\\
	\vspace{0.1cm}
	\noindent\large\textbf{Campion School Bhopal} \hfill April 2005 -- March 2017 \\ Primary and Middle School \hfill CGPA: 10/10\\
\end{flushleft}

% -----------------------------------

\section*{Internships}

\noindent\large\textbf{Web and Mobile app development Intern at C.R.I.S. Secunderabad} \hfill 1 June 2025 -- 30 July 2025

\vspace{2pt}
\noindent\hspace{15pt}{\fontsize{11.5pt}{13.5pt}\selectfont \textbf{Keywords:} Git, Flutter, Dart, Javascript, Nodejs, MySQL, POSTMAN} \hfill \href{https://gitlab.com/rishabhsahu325/cris_sec25_divjan}{\faIcon{link}\textbf{Project repository}}

\begin{itemize}[leftmargin=15pt, topsep=0pt, itemsep=2pt]
	\item {\fontsize{11.5pt}{13.5pt}\selectfont Worked on a prototype mobile application for Divyangjan card application website}
	\item {\fontsize{11.5pt}{13.5pt}\selectfont Implemented backend API routes}
	\item {\fontsize{11.5pt}{13.5pt}\selectfont Contributed in database design and testing}
	\item {\fontsize{11.5pt}{13.5pt}\selectfont Added frontend pages for Railway division wise reports}
\end{itemize}


% -----------------------------------

\section*{Projects}

% --- PROJECT 1: JSON API SERVER ---
\begin{comment}
description:
A small backend service exposing teacher and course data from a
Postgres database as JSON. Built as a sandbox for Flask routing
patterns, request handling, and one-command containerised deployment
with Docker and Bash.
\end{comment}
\subsection*{JSON API server \hfill {\normalfont\itshape\small Backend Development}}

\vspace{2pt}
\noindent\hspace{15pt}{\fontsize{11.5pt}{13.5pt}\selectfont \textbf{Keywords:} Flask, Docker, Python, API design, Postgresql \hfill \href{https://gitlab.com/projectsa2/jsonapiserver}{\faIcon{link}\textbf{Project repository}}}

\begin{itemize}[leftmargin=15pt, topsep=0pt, itemsep=2pt]
	\item {\fontsize{11.5pt}{13.5pt}\selectfont Handles external http requests to fetch data from Postgresql Database}
	\item {\fontsize{11.5pt}{13.5pt}\selectfont Perform create and view operations for teacher and course data in the database}
	\item {\fontsize{11.5pt}{13.5pt}\selectfont Retrieve responses to http requests in JSON format}
	\item {\fontsize{11.5pt}{13.5pt}\selectfont Utilized Docker and BASH to reduce setup time}
\end{itemize}


\centerline{\rule{0.8\textwidth}{0.05pt}}

% --- PROJECT : Wine classifier based on decision tree  ---
\begin{comment}
description:
A decision-tree classifier trained on sklearn's wine dataset. Used to
walk through the full supervised-learning workflow end to end -- data
loading, model fitting with scikit-learn, and evaluation through
Pyspark's model evaluator.
\end{comment}
\subsection*{Wine Classifier based on Decision Tree \hfill {\normalfont\itshape\small Machine Learning}}

\noindent\hspace{15pt}{\fontsize{11.5pt}{13.5pt}\selectfont \textbf{Keywords:} Python, Pyspark, Pandas, Decision Tree classification \hfill \href{https://gitlab.com/projectsa2/wineclassifier}{\faIcon{link}\textbf{Project repository}}}

\begin{itemize}[leftmargin=15pt, topsep=0pt, itemsep=2pt]
	\item {\fontsize{11.5pt}{13.5pt}\selectfont Generated ML classifier model to classify Wine from characteristics}
	\item {\fontsize{11.5pt}{13.5pt}\selectfont Classifier trained on sklearn's wine dataset}
	\item {\fontsize{11.5pt}{13.5pt}\selectfont Evaluated trained model through pyspark's model evaluator}
	\item {\fontsize{11.5pt}{13.5pt}\selectfont Achieved reasonably good accuracy of 85\% on test metrics}

\end{itemize}

\centerline{\rule{0.8\textwidth}{0.05pt}}

% --- PROJECT : N-gram Language Model ---
\begin{comment}
description:
A typing word predictor that reads the starting letters and fills the complete word by choosing among prediction options generated by n gram model trained on a text corpus.
\end{comment}
\subsection*{N-gram Language Model for Next-Word Prediction\hfill {\normalfont\itshape\small Machine Learning}}

\noindent\hspace{15pt}{\fontsize{11.5pt}{13.5pt}\selectfont \textbf{Keywords:} Python, Character N-gram, Curses (TUI)  \hfill \href{https://gitlab.com/projectsa2/wineclassifier}{\faIcon{link}\textbf{Project repository}}}

\begin{itemize}[leftmargin=15pt, topsep=0pt, itemsep=2pt]
	\item {\itemsize Developed a Character-level N-gram model to suggest word completions in real-time as the user types.}
	\item {\itemsize Implemented Maximum Likelihood Estimation (MLE) to calculate transition probabilities between character sequences.}
	\item {\itemsize Built a low-latency Terminal User Interface (TUI) using \texttt{curses} that tracks keystroke efficiency metrics live.}
\end{itemize}

\centerline{\rule{0.8\textwidth}{0.05pt}}

% --- PROJECT : Secure Telemedical ---
\begin{comment}
description:
This project implements a simulation of a secure telemedical conferencing system that enables confidential communication between a doctor and multiple patients.
The system uses cryptographic techniques including ElGamal encryption/digital signatures and AES encryption to ensure secure authentication, key exchange, and encrypted communication.
\end{comment}
\subsection*{Secure many to one communication system with Digital Signatures\hfill {\normalfont\itshape\small Cybersecurity}}
\noindent\hspace{15pt}{\fontsize{11.5pt}{13.5pt}\selectfont \textbf{Keywords:} Python, Cryptography, Network Security, Socket Programming  \hfill \href{https://gitlab.com/iiith2426-ris/assignments/sns/a2-telemedical-ds}{\faIcon{link}\textbf{Project repository}}}

\begin{itemize}[leftmargin=15pt, topsep=0pt, itemsep=2pt]
	\item {\itemsize Engineered a secure communication protocol ensuring confidentiality using AES-256 and ElGamal encryption}
	\item {\itemsize Implemented mutual authentication and integrity verification mechanisms using ElGamal Digital Signatures}
	\item {\itemsize Mitigated replay attacks via timestamp validation and enforced access control through temporary blocking}
	\item {\itemsize Orchestrated secure session management with dynamic group key generation for encrypted broadcast messaging}
\end{itemize}

\centerline{\rule{0.8\textwidth}{0.05pt}}

% --- PROJECT : Multithreaded HTTP Server ---
\begin{comment}
description:
This project implements a simple multithreaded server in rust that caters to multiple http requests through a thread pool and handles concurrency through Rust primitives like Arc, Mutex and channels.
\end{comment}
\subsection*{Multithreaded HTTP Server\hfill {\normalfont\itshape\small Systems Programming}}
\noindent\hspace{15pt}{\fontsize{11.5pt}{13.5pt}\selectfont \textbf{Keywords:} Rust, Multithreading, TCP, HTTP, Thread Pool, Concurrency  \hfill \href{https://github.com/RishabhSahuIIIT/rustMultithreadServer}{\faIcon{link}\textbf{Project repository}}}
\begin{itemize}[leftmargin=15pt, topsep=0pt, itemsep=2pt]
	\item {\fontsize{11.5pt}{13.5pt}\selectfont Built an HTTP/1.1 server from scratch using Rust's standard library }
	\item {\fontsize{11.5pt}{13.5pt}\selectfont Implemented a fixed-size thread pool dispatching jobs across workers}
	\item {\fontsize{11.5pt}{13.5pt}\selectfont Handled GET request routing with \texttt{200 OK}, simulated latency \texttt{/sleep}, and \texttt{404 NOT FOUND} responses}
	\item {\fontsize{11.5pt}{13.5pt}\selectfont Applied Rust ownership and concurrency primitives like Arc, Mutex and channels for safe parallel execution}
\end{itemize}

\centerline{\rule{0.8\textwidth}{0.05pt}}

% --- PROJECT : Custom Shell ---
\begin{comment}
description:
A custom C++ shell that replicates core command-line functionalities. Engineered to handle external program execution, output redirection, and inter-process communication using native Linux APIs.
\end{comment}
\subsection*{Shell \hfill {\normalfont\itshape\small Systems Programming}}

\noindent\hspace{15pt}{\fontsize{11.5pt}{13.5pt}\selectfont \textbf{Keywords:} C++, Linux API, Process handling, System calls \hfill \href{https://github.com/RishabhSahuIIIT/codecrafters-shell-cpp}{\faIcon{link}\textbf{Project repository}}}

\begin{itemize}[leftmargin=15pt, topsep=0pt, itemsep=2pt]
	\item {\fontsize{11.5pt}{13.5pt}\selectfont Implemented external program execution through fork and exec}
	\item {\fontsize{11.5pt}{13.5pt}\selectfont Developed shell built-ins including \texttt{cd}, utilizing OS function calls for directory changes rather than internal state management}
	\item {\fontsize{11.5pt}{13.5pt}\selectfont Integrated the readline library for tab-autocompletion, applied specifically to built-in commands}
	\item {\fontsize{11.5pt}{13.5pt}\selectfont Orchestrated inter-process communication via pipes}
	\item {\fontsize{11.5pt}{13.5pt}\selectfont Manipulated file descriptors to enable output redirection}
\end{itemize}

\centerline{\rule{0.8\textwidth}{0.05pt}}

% --- PROJECT : HTML Autocomplete with Lex ---
\begin{comment}
description:
A command-line tool that performs lexical analysis on HTML files using GNU Flex
and provides autocomplete suggestions for HTML tags and attributes. A Trie data
structure built in C matches partial token prefixes against all recognised HTML
tokens, demonstrating practical application of compiler front-end techniques.
\end{comment}
\subsection*{HTML Autocomplete using Lexer \hfill {\normalfont\itshape\small Compilers}}

\noindent\hspace{15pt}{\fontsize{11.5pt}{13.5pt}\selectfont \textbf{Keywords:} GNU Flex/Lex, C, Lexical Analysis, Trie, Compiler Design \hfill \href{https://gitlab.com/iiith2426-ris/assignments/compilers/a2-html-autocomplete-lex}{\faIcon{link}\textbf{Project repository}}}

\begin{itemize}[leftmargin=15pt, topsep=0pt, itemsep=2pt]
	\item {\itemsize Implemented an HTML lexer using GNU Flex with custom grammar rules to tokenise HTML documents}
	\item {\itemsize Built a Trie-based string matching engine in C to provide autocomplete suggestions from partial HTML tags and attributes}
	\item {\itemsize Extended token coverage beyond standard HTML spec to support tags like \texttt{img}, \texttt{header}, and \texttt{srcset}}
\end{itemize}

\centerline{\rule{0.8\textwidth}{0.05pt}}

% --- PROJECT : Network Intrusion Detection and Prevention System ---
\begin{comment}
description:
A Network-based Intrusion Detection and Prevention System (NIDPS) that monitors
live TCP traffic using Scapy, detects port scanning and OS fingerprinting attacks
via sliding-window signature analysis, and automatically blocks malicious IPs
using platform-specific firewall commands. Multi-threaded design keeps the CLI
responsive while capture runs in the background.
\end{comment}
\subsection*{Network Intrusion Detection and Prevention System \hfill {\normalfont\itshape\small Cybersecurity}}

\noindent\hspace{15pt}{\fontsize{11.5pt}{13.5pt}\selectfont \textbf{Keywords:} Python, Scapy, Network Security, iptables, Multithreading \hfill \href{https://gitlab.com/iiith2426-ris/assignments/sns/a3}{\faIcon{link}\textbf{Project repository}}}

\begin{itemize}[leftmargin=15pt, topsep=0pt, itemsep=2pt]
	\item {\itemsize Performed real-time TCP packet capture and analysis using Scapy's low-level network sniffing capabilities}
	\item {\itemsize Implemented signature-based detection for port scanning (sliding time-window) and OS fingerprinting (TCP flag pattern matching)}
	\item {\itemsize Automated threat response by dynamically blocking attacker IPs via iptables (Linux) and Windows Firewall (Windows)}
	\item {\itemsize Decoupled packet capture from the management CLI using multi-threading for continuous monitoring without UI blocking}
\end{itemize}

\centerline{\rule{0.8\textwidth}{0.05pt}}

% --- PROJECT : Agentic AI Research Assistant ---
\begin{comment}
description:
A multi-agent system built on the Moya framework that processes academic PDF
papers end-to-end: parsing with PyMuPDF, structured summarisation, cross-paper
synthesis, and mini-survey generation with inline citations. A dual Ollama
architecture runs orchestration and task execution on separate local LLM instances
to prevent blocking, with every decision and output logged to timestamped JSONL
trace files for full observability.
\end{comment}
\subsection*{Agentic AI Research Assistant \hfill {\normalfont\itshape\small Artificial Intelligence}}

\noindent\hspace{15pt}{\fontsize{11.5pt}{13.5pt}\selectfont \textbf{Keywords:} Python, Moya framework, Ollama, Llama 3.1, PyMuPDF, Multi-agent systems \hfill \href{https://github.com/RishabhSahuIIIT/AgenticAIResearchAssistant_Moya}{\faIcon{link}\textbf{Project repository}}}

\begin{itemize}[leftmargin=15pt, topsep=0pt, itemsep=2pt]
	\item {\itemsize Built a multi-agent pipeline using the Moya framework to parse, summarise, synthesise, and survey academic PDF papers}
	\item {\itemsize Orchestrated dual Ollama (Llama 3.1) instances on separate ports to decouple decision-making from task execution, preventing LLM blocking}
	\item {\itemsize Implemented full observability by logging every agent decision, LLM prompt, and output to timestamped JSONL trace files per run}
	\item {\itemsize Generated structured mini-surveys ($\leq$800 words) with inline citations synthesised across multiple research papers}
\end{itemize}

\centerline{\rule{0.8\textwidth}{0.05pt}}

% --- PROJECT : Citestat ---
\begin{comment}
description:
A browser-only academic citation analytics dashboard that fetches paper and author
metadata from the Crossref and OpenCitations REST APIs. Users can look up
any DOI or ORCID to see yearwise citation histograms (Chart.js), citation network
graphs (Reagraph / D3.js), and bibliometric impact metrics (h-index, m-quotient)
computed entirely on the client. Deployed on Vercel with automatic CI/CD on each
merge to master.
\end{comment}
\subsection*{Citestat -- Academic Citation Analytics Dashboard \hfill {\normalfont\itshape\small Web Development}}

\noindent\hspace{15pt}{\fontsize{11.5pt}{13.5pt}\selectfont \textbf{Keywords:} React, TypeScript, Vite, Tailwind CSS, Chart.js, Reagraph, Crossref API, OpenCitations API \hfill \href{https://github.com/RishabhSahuIIIT/citestat}{\faIcon{link}\textbf{Project repository}}}

\begin{itemize}[leftmargin=15pt, topsep=0pt, itemsep=2pt]
	\item {\itemsize Built a client-side citation analytics dashboard that queries Crossref and OpenCitations REST APIs for paper and author metadata}
	\item {\itemsize Visualised yearly citation trends as histograms using Chart.js and citation network graphs using Reagraph (backed by D3.js)}
	\item {\itemsize Computed bibliometric impact metrics (h-index, m-quotient) fully in the browser with no backend dependency}
	\item {\itemsize Deployed to Vercel's free tier with CI/CD; every merge to master automatically triggers a production deployment at \href{https://citestat.vercel.app}{citestat.vercel.app}}
\end{itemize}

\centerline{\rule{0.8\textwidth}{0.05pt}}

% --- PROJECT : CRDT Collaborative Text Editor ---
\begin{comment}
description:
A fully-distributed real-time collaborative text editor built on two complementary
CRDT approaches: an operation-based Replicated Growable Array (RGA) for document
text and a state-based Last-Write-Wins (LWW) Register for metadata. Multiple
users can concurrently edit the same document over a WebSocket relay with
guaranteed convergence — even after network partitions or reordered delivery —
plus offline editing backed by IndexedDB, HMAC-based access control, and cursor
presence.
\end{comment}
\subsection*{Real-Time Collaborative Text Editor using CRDT \hfill {\normalfont\itshape\small Distributed Systems}}

\noindent\hspace{15pt}{\fontsize{11.5pt}{13.5pt}\selectfont \textbf{Keywords:} TypeScript, React, CodeMirror 6, WebSocket, CRDT, RGA, LWW Register, IndexedDB \hfill \href{https://github.com/RishabhSahuIIIT/CRDT_collaborative_text_editor}{\faIcon{link}\textbf{Project repository}}}

\begin{itemize}[leftmargin=15pt, topsep=0pt, itemsep=2pt]
	\item {\itemsize Implemented an operation-based RGA CRDT for document text, ensuring convergence across concurrent edits via vector clocks and causal delivery}
	\item {\itemsize Implemented a state-based LWW Register CRDT for metadata (title, cursors), with timestamp + siteId total order for deterministic conflict resolution}
	\item {\itemsize Built a WebSocket relay server and React/CodeMirror 6 browser editor with IndexedDB-backed offline editing and automatic reconnect sync}
	\item {\itemsize Enforced per-message access control using HMAC-signed role tokens, applied at the relay so role changes take effect immediately without reconnection}
\end{itemize}

% -----------------------------------
\section*{Certifications}
\begin{itemize}
	\item \href{https://www.educative.io/verify-certificate/Y5xyKBov0ppTp41gXrqlx5c3LGP7RAWBVHJ}{\faIcon{link}Working with containers}: Docker, Docker compose, Docker Swarm (EDUCATIVE.IO)
	\item \href{https://www.educative.io/verify-certificate/j2l3BzfZ3O94zqGqKSjXR9PNLqr7uA}{\faIcon{link}Flask Develop web applications in Python}: Flask, Python, CSS, HTML (EDUCATIVE.IO)
	\item \href{https://www.hackerrank.com/certificates/388867dc0d5f}{\faIcon{link}Java basic certification} (HACKERRANK)
\end{itemize}

% -----------------------------------
\section*{Skills}
% Using a simple itemize with zero margin allows text to flow naturally
\begin{itemize}[leftmargin=*, itemsep=0pt, parsep=0pt]
	\item \textbf{Programming and Databases :} Python, C/C++, Rust, Java, SQL, MongoDB
	\item \textbf{Web Technologies:} Flask, React, Express, HTML5, CSS, Javascript
	\item \textbf{DevOps \& Tools:} BASH, Linux Shell, Docker, Git
	\item \textbf{Core Concepts:} Data Structures and Algorithms, Object Oriented Programming
\end{itemize}

% -----------------------------------
\section*{Accomplishments}
\begin{itemize}
	\item Accomplished 77 All India Rank in IIIT PGEE CSE 2024 exam
	\item Achieved 95.4 percentile in GATE exam CS/IT 2024
	\item Qualified Stage 1 of N.T.S.E. in Madhya Pradesh
	\item Zonal Gold medalist in SOF National cyber Olympiad Madhya Pradesh zone class 9th
\end{itemize}

% -----------------------------------
\section*{Interests and Extracurricular Activities}
\begin{itemize}
	\item Playing badminton
	\item Participated in Discernment International Summit M.U.N. 2016-17 at Noida
\end{itemize}
\end{document}
